\documentclass[10pt,a4paper]{amsart}
\usepackage[margin=19mm]{geometry}
\usepackage{amsmath,amssymb,mathtools}
\usepackage[scr=rsfs,cal=euler]{mathalfa}
\usepackage{xfrac}
\usepackage[unicode,hidelinks,bookmarks=false]{hyperref}
\newcommand{\ie}{i.e.\ }
\newcommand{\cf}{cf.\ }
\newcommand{\resp}{respectively\ }
\newcommand{\Subsection}{\subsection}
\providecommand{\nobreakdash}{\nobreak\hbox{-}}
\setlength{\emergencystretch}{2em}
\multlinegap0pt
\usepackage{amssymb}
\usepackage{dsfont}
\usepackage{bm}
\newtheorem{theorem}{Theorem}
\DeclareMathOperator{\CF}{CF}
\DeclareMathOperator{\Fix}{Fix}
\DeclareMathOperator{\Gr}{Gr}
\DeclareMathOperator{\HF}{HF}
\newcommand{\R}{{\mathbb{R}}}
\DeclareMathOperator{\Symp}{Symp}
\newcommand{\p}{\partial}
\DeclareMathOperator{\supp}{supp}
\title{Floer-theoretic entropy of exact symplectomorphisms}
\author{Joontae Kim, Myeonggi Kwon}
\date{Source: arXiv:2608.21780v2 (CC BY 4.0)}
\begin{document}
\maketitle

\noindent\emph{Source and licence.} Floer-theoretic entropy of exact symplectomorphisms. Version arXiv:2608.21780v2 (CC BY 4.0), distributed by arXiv under the Creative Commons Attribution 4.0 International licence, http://creativecommons.org/licenses/by/4.0/, as recorded in the arXiv licence metadata for this exact version and shown on the abstract page https://arxiv.org/abs/2608.21780v2. Full licence notice: \texttt{licenses/arxiv-2608.21780-CC-BY-4.0.txt}.\par\smallskip

\noindent\textit{Editorial scope.} Counted unit: the article's theorem can\_isom (source0.tex lines 477--486). The article's complete original proof of it is delivered with it (source0.tex lines 487--500, one \texttt{proof} environment of 14 lines). No mathematics is written, repaired or re-derived: the statement and the proof are the article's own lines, in the article's own notation, and no retained line is substituted or re-flowed.  No further source-local context is needed: every cross-reference of every retained line resolves inside the two delivered spans.  Nothing was omitted from any retained span, so the package carries no omission record.  No retained line contains a citation, so the wrapper carries no bibliography and no bibliography entry.  No retained line contains a figure environment, an image-inclusion command or drawing code, so no image is delivered and the package declares no asset dependency.  Editorial formatting only, in addition: an AMS \texttt{amsart} preamble that restates the article's own declarations and macros used by the retained lines, namely the environments \texttt{theorem} on separate counters, together with the standard packages those lines need.  The retained environments print in this document as Theorem 1 in order of appearance, whereas the article prints the same items with the numbering of its own full document.  The complete native source of the article as distributed in the arXiv e-print for this exact version is bundled unchanged as \texttt{sources/208276/source0.tex}.
\begin{theorem}\label{thm: can_isom}
%Given regular ${\bm J}_\phi$, there exists a canonical identification between two chain complexes
%$$
%(CF(\phi),\p^{{\bm J}_\phi}) \cong (CF(\Delta,\Gr(\phi)),\p^{{\bm J}_X}).
%$$
For every $\phi\in\Symp(W,\lambda)$ there exists a canonical isomorphism
$$
\HF(\phi)\cong \HF(\Delta,\Gr(\phi)).
$$
\end{theorem}

\begin{proof}
By the invariance properties, we can assume that $\phi$ is non-degenerate, equivalently, $\Delta$ and $\Gr(\phi)$ intersect transversely.
Since the generators of $\CF(\phi)$ and $\CF(\Delta,\Gr(\phi))$ are in bijection under $\Fix(\phi)\ni x\mapsto (x,x)\in \Delta\cap\Gr(\phi)$, it suffices to show that the differentials coincide.
For generic ${\bm J}=\{J_t\}_{t\in \R}\in\mathcal{J}_{\phi}$ we consider the family $\widetilde{{\bm J}}=\big\{(-J_{\frac{1-t}{2}})\oplus J_{\frac{1+t}{2}}\big\}_{t\in[0,1]}$ of split almost complex structures on $\widetilde{W\times W}$.
The moduli space of twisted Floer trajectories for ${\bm J}$ corresponds bijectively to that of Floer strips for $\widetilde{{\bm J}}$ via
$
u(s,t) \mapsto\widetilde{u}(s,t)=\left(u\big(\frac{s}{2},\frac{1-t}{2}\big),u\big(\frac{s}{2},\frac{1+t}{2}\big)\right)
$.
Moreover, if $u$ is regular, then  so is $\widetilde{u}$, where regularity means that the associated Fredholm operator is surjective.
A priori, $\widetilde{u}$ is contained in $W\times W$, but the argument below shows that it actually lies in $\widetilde{W\times W}$.
Using the assumption $\supp(\phi)\subset W\setminus \p W$ and the maximum principle, we can choose a compact region $K\subset W\setminus\p W$ that contains all fixed points of $\phi$ and all twisted Floer trajectories contributing to $\HF(\phi)$.
We assume that the hypersurface $\Sigma\subset W\times W$ is chosen to be disjoint from $K\times K$, while still satisfying all of our standing assumptions.
With this choice, the correspondence above ensures that every $\widetilde{u}$ lies in $\widetilde{W\times W}$, which finishes the proof.
\end{proof}

\end{document}
